%!TEX root = main.tex

Motivated by an old problem in additive combinatorics about estimating the size of Sidon sets \cite{CRV}, many authors have studied the problem of finding the best constant $c_{\max}$ such that for any function $f\in L^{1}(\R)$ supported in $[-1/4,1/4]$ the following inequality holds
$$
\text{max}_{-1/2 \le t \le 1/2} \int_{\R} f(t-x)f(x) \, dx \ge c_{\max} \left(\int_{-1/4}^{1/4} f(x) \, dx\right)^2.
$$
The best known result so far was obtained by Cloninger and Steinerberger in \cite{CS}, they proved that $c \ge 1,28.$ Many other lower bounds were previously obtained in \cite{CRT,G,MOB1,Y,MOB2,MV}.
Inspired by this question, two other related problems were proposed and studied by Barnard and Steinerberger in \cite{BarnardSteinerberger}. One of their results was the following: The inequality
\begin{equation}\label{eq:l1l2}
\int_{-1/2}^{1/2}\int_{\R}f(x)f(x+t) dx dt\leq 0.91\|f\|_1\|f\|_2,
\end{equation}
holds for any function $f \in L^1(\R) \cap L^2(\R)$. It was also established in \cite{BarnardSteinerberger} via an example, that the best constant such that \eqref{eq:l1l2} holds is at least 0.8.
The upper bound was recently improved by the second author
and Ramos in \cite{MR}, where they proved that this inequality still hold true when we write 0.865 instead of 0.91. A natural question is, what happen when we consider a different probability space, in particular, what happen when we consider Gaussian means. This question was also addressed in \cite{MR}:
\begin{proposition}[\cite{MR}, Theorem 1.2]\label{Thm JJ}
Let $a$ be a positive real number. For any $f\in L^{1}(\R)\cap L^{2}(\R)$. The following inequality holds
\begin{align}\label{MR thm 1}
\left(\frac{a}{\pi}\right)^{1/2}\int_{\R}\int_{\R}f(x)f(x+t)e^{-a t^{2}}dxdt\leq \left(\frac{8a}{27\pi}\right)^{1/4}\|f\|_1\|f\|_2,
\end{align}
and $\left(\frac{8a}{27\pi}\right)^{1/4}$ can not be replaced by $\left(\frac{a}{4\pi}\right)^{1/4}$. 
\end{proposition}
%In particular, if $a=2\pi$ the upper bound is 0.8773 and the lower bound 0.8408.$.

	%{\it{Structure of the paper:}} 
	In this paper, using Euler-Lagrange equations, in Section 2 we establish the existence of extremizers for \eqref{eq:l1l2}, \eqref{MR thm 1} and a more general class of weights (nonnegative functions $w:\mathbb{R}\to[0,+\infty)$). In Section 3, via a discretization argument, we find an almost optimal numerical approximation for the best constants allowed in these inequalities, see Table 1. Finally, in Section 4, we discuss a related problem about autoconvolution.\\ %Finally, in the appendix, we explain how some extra regularity assumptions on the weight (for instance when dealing with Gaussian functions) allow us to obtain even better approximations (as expected).\\
	
	Related results about autoconvolution inequalities were recently obtained in \cite{CJLL}. We hope variations of the methods outlined in this paper can be applied to such, and other, related problems. As another example,  a classical problem proposed by Erdos, the Minimum overlaping problem, can be reformulated in terms of autorrelations as observed by Haugland in \cite{Ha}, we believe that there is a strong relation between this problem and the one analyzed in this manuscript.

\newcommand{\gr}[1]{{\color{red}#1}}
\begin{table}
\renewcommand{\arraystretch}{1.5}
	\begin{tabular}{ccccc}
		\textbf{Weight} & \multicolumn{3}{c}{\textbf{Spectral Method}} & \textbf{Fixed Point}  \\
			 					& {\small Lower bound} & {\small Upper bound}  & {\small Difference} 	 	& {\small Lower bound} 

		\\
		\hline

		$\chi_{[-1/2,1/2]}$ & $0.80558\gr{09}$ & $0.80558\gr{96}$ & $<9\cdot 10^{-6}$ & $0.80558\gr{09}$\\
		$\exp(-\pi x^2)$    & $0.7152\gr{474}$ & $0.7152\gr{576}$ & $<1.2\cdot 10^{-5}$ & $0.7152\gr{475}$
	\end{tabular}

	\caption{Upper and Lower bounds for the average problem. These bounds have been found with the algorithm described in Section 3.3, with a value of $\delta \approx 1.45\cdot 10^{-3}$ and $\Delta \lambda \approx 0.001$. The implementation can be found at \cite{Github}.}
	\label{table}
\end{table}
	

\begin{figure}
    \centering
    \includegraphics[width=.9\textwidth]{extremizers.png}
    \caption{Numerical extremizers for the discretized problem. Extremizers are normalized so that $\|f^*\|_1\|f^*\|_2 = 1$. }
    \label{fig:my_label}
\end{figure}